%% ================================================================
%%  III. HOMOGENEOUS STABILITY DIAGNOSTIC
%% ================================================================
\section{Homogeneous Stability Diagnostic}
\label{sec:lsa}

\subsection{Lumped diagnostic and base state}

We use a homogeneous/lumped linear diagnostic for spatially uniform
states of the PDE system~\eqref{eq:governing}.  This diagnostic is
not the full spectrum of the linearized Robin-boundary PDE. It
combines a lumped homogeneous balance, including boundary exchange
terms, with local bulk dispersion corrections. It is used as a
necessary but not sufficient indicator of possible oscillatory
feedback.  The uniform steady states $(J_0,u_0,\theta_0)$ satisfy
three algebraic balances:
\begin{align}
  \mu(J_0,\theta_0)&=\mu_b, &
  \Bi_c(1-u_0)&=\Da\,J_0\,R_{b0}, \nonumber\\
  \Bi_T\theta_0&=\Da\,J_0\,R_{b0},
  \label{eq:lsa_ss}
\end{align}
where $R_{b0}=R_b(u_0,J_0,\theta_0)$.  Linearizing the lumped
balance about this base state and adding local bulk-gradient
corrections gives a $3\times 3$ eigenvalue problem for each
diagnostic wavenumber $k$ (Appendix~\ref{app:lsa}):
\begin{equation}
  \sigma\,\widehat{\mathbf y} = \mathbf A(k)\,\widehat{\mathbf y},
  \qquad
  \mathbf A(k)=\mathbf A_0 + k^2\mathbf D_2 + k^4\mathbf D_4.
  \label{eq:dispersion_matrix}
\end{equation}
Here $\mathbf A_0$ encodes reaction kinetics and surface exchange,
$\mathbf D_2$ collects diffusive transport, and $\mathbf D_4$ arises
from the Cahn--Hilliard gradient-energy regularization.  Uniform
branches are classified as \emph{stable}
($\Re\sigma_1(0)<0$), \emph{Hopf unstable}
($\Re\sigma_1(0)>0$, $\Im\sigma_1(0)\neq 0$), or
\emph{monotone unstable} ($\Re\sigma_1(0)>0$,
$\Im\sigma_1(0)=0$).

\subsection{Two instability mechanisms}

The dispersion relation reveals two physically distinct instabilities
(Fig.~\ref{fig:stability}).

\paragraph{Hopf instability ($k=0$).}
The long-wave mode is governed by $\mathbf A_0$ alone.  The positive
feedback originates in the heat equation: when the Arrhenius
heat--reaction coupling $\Da\,J_0 R_\theta$ exceeds the surface heat
loss $\Bi_T$, a temperature perturbation amplifies the reaction rate,
which in turn raises the temperature further.
The negative feedback follows the off-diagonal loop
$\theta\to J\to u\to\theta$: rising temperature increases $\chi$ and
drives collapse, collapse suppresses reactant supply, and reduced
reactant weakens the heat source.  Oscillation requires these two
feedback paths to operate on well-separated timescales; the onset
condition is set by the balance of the characteristic-polynomial
coefficients of $\mathbf A_0$
(Appendix~\ref{app:lsa}).

\paragraph{Spinodal instability ($k>0$).}
When the thermodynamic derivative $f_J\equiv\partial_J[\,m_\mathrm{mix}
+m_\mathrm{el}\,]<0$, the base state lies inside the Flory--Huggins
spinodal region and the $(1,1)$ entry of $\mathbf D_2$ acts as a
negative diffusivity: spatial perturbations in $J$ are amplified
rather than damped~\cite{Onuki1993,Bai2023}.  The $k^4$ term limits
growth to a finite band with maximum at
$k^\ast\approx\sqrt{-f_J/(2\ell^2)}$
[Fig.~\ref{fig:stability}(c)].  When triggered in isolation
($\Da=0$, fixed $\theta$ above LCST), this mode produces alternating
swollen and collapsed domains (Appendix~\ref{app:spinodal}).

\paragraph{Why spinodal does not manifest during oscillation.}
Although the spinodal growth rate ($\sigma^\ast\sim 10^3$) far exceeds
the Hopf rate ($\sim\!1$), three factors prevent spinodal patterns
from developing during the oscillation cycle: (i)~the system traverses
the spinodal region only transiently during collapse
($\lesssim 0.5\,\tau$, as confirmed in Sec.~\ref{sec:results}), too briefly for short-wavelength domains to
reach macroscopic amplitude; (ii)~the Thiele penetration gradient and
Robin boundary conditions break the translational symmetry required
for classical spinodal pattern selection; (iii)~the large-amplitude
Hopf limit cycle continuously homogenizes the spatial structure at
each period.

\begin{figure*}[tbp]
  \centering
  \includegraphics[width=\textwidth]{fig1.pdf}
  \caption{Linear stability of the 1D slab.
    \textbf{(a)}~Steady-state branches $J_0$ versus $\Bi_T$,
    colored by $k=0$ stability; the star marks the working point.
    \textbf{(b)}~Leading eigenvalue (real and imaginary parts) along
    the collapsed branch, showing the monotone-to-Hopf transition.
    \textbf{(c)}~Full dispersion curves for three $\Bi_T$ values;
    inset resolves the $k\to 0$ Hopf regime.
    \textbf{(d,e)}~Stability maps in the $(\Da,\Bi_T)$ and
    $(S_\chi,\Bi_T)$ planes.
    \textbf{(f)}~Linear period $T_\mathrm{LSA}=2\pi/\omega$ and
    growth rate vs $\Da$.}
  \label{fig:stability}
\end{figure*}

\subsection{Implications for the working point}
\label{sec:lsa_implications}

The lower-$J$ (collapsed) branch at the default parameters undergoes
a sequence
monotone~$\to$~Hopf~$\to$~stable as $\Bi_T$ increases: weak cooling
causes thermal runaway, intermediate cooling produces oscillatory
growth, and strong cooling stabilizes the system
[Fig.~\ref{fig:stability}(a,b)].  The Hopf region persists only when
$\Da$ and $S_\chi$ are large enough to overcome surface heat loss
[Fig.~\ref{fig:stability}(d,e)].  The linear period estimate
$T_\mathrm{LSA}=2\pi/\omega$ [Fig.~\ref{fig:stability}(f)] provides
the small-amplitude precursor to the nonlinear oscillation period; the
dispersion analysis correctly identifies the existence, type, and
timescale of the instability that seeds the slab dynamics.


